% Propietario conservado: /Users/ruben/Documents/New project/output/EXTREMA_MEDIA_RAZON_RECIPROCIDAD_ES_EN_20260918/ARTICULO/ES/source/sections/generacion.tex
\section{Coinductive generation of \texorpdfstring{\(\pi,\varphi,e\)}{pi, phi, e} to arbitrary depth}
\label{x_reciprocidad:sec:generacion}

Emission of six blocks is a finite stage of the construction, not a complete decimal expansion. Its function is to produce regions with multiplicity, orientation and transport rules. Prolongation must preserve these data and determine a compatible family of prefixes of arbitrary length. The word \emph{generation} denotes this operation; \emph{genealogy} denotes reconstruction of its dependencies. Both notions enter into the process but serve distinct functions.

\subsection{Orbits, regions, and characters of closure, propagation, and self-scaling}

On a word \(w=(w_1,\ldots,w_6)\), define
\[
 r(w)=(w_3,w_1,w_2,w_5,w_6,w_4),\qquad
 s(w)=(w_4,w_5,w_6,w_1,w_2,w_3).
\]
We have \(r^3=s^2=I\) and \(srs=r^{-1}\); these permutations therefore realize the dihedral group \(D_3\). The action preserves the catalogue and multiplicities and commutes with \(\Psi\). The quotient of the \(243\) visible words has \(43\) orbits: thirty-eight of cardinality six and five of cardinality three.

The closure class is recognized by a fiber predicate: its six orientations have five emissions per word, total multiplicity \(1008\) and distribution \(432+4\cdot144\). There is only one orbit with these properties in the catalogue. To make the other two selections precise, write \(f(w)=\#\Psi^{-1}(w)\), \(m(w)=\sum_{U\in\Psi^{-1}(w)}\operatorname{mult}(U)\) and \(\nu(w)=(n_0,n_1,n_2)\). The additive diagonal class of an emission is \(d(U)=[i+j]_9\) in indices from zero to eight; its additive signature preserves this class and the orientation \(ES\) or \(NO\). Let \(\mathcal D(\mathcal O)\) be the set of these classes for all emissions over an orbit.

Among the orbits of size six, the joint predicate
\[
 f(w)=1,\qquad m(w)=144\quad(w\in\mathcal O),\qquad
 \mathcal D(\mathcal O)=\{0,3,6\}
\]
leaves exactly six orbits. In this sector, the census \(\nu=(2,3,1)\), equal to that of closure, selects only propagation; the census \(\nu=(2,2,2)\) selects only self-scale. The diagonal support cannot be omitted: without it there are four balanced single-point orbits of multiplicity \(144\), not one. Finally, existence of an emission with \(d=6\) and orientation \(ES\) selects one representative of each orbit and leads to
\begin{equation}
 w^{\rm cl}=010211,\qquad w^{\rm pr}=201101,\qquad w^{\rm au}=121200.
 \label{x_reciprocidad:eq:palabras-regionales}
\end{equation}
These symbols first denote catalogue regions. Scalar identification will be proved through their readers; they are not chosen because they match digits of a known constant.

The closure word \(010211\) has the following five preimages:
\begin{center}\small
\begin{tabular}{@{}lll r@{}}\toprule
Six-block emission&Firma multiplicativa&Role&Multiplicity\\\midrule
\(501,614,498,169,272,272\)&\(555555\)&transverse&432\\
\(810,923,870,169,272,272\)&\(339555\)&positive orientation&144\\
\(810,983,810,169,272,272\)&\(393555\)&positive orientation&144\\
\(870,923,810,169,272,272\)&\(933555\)&positive orientation&144\\
\(870,923,810,418,674,521\)&\(933717\)&oriented return&144\\\bottomrule
\end{tabular}
\end{center}
The multiplicity \(432\) and constant signature \(555555\) characterize the transverse region. They do not identify the central position \((5,5)\) of APP. Confusing a six-window signature with a support position would remove the distinction between numerical region and electronic structure.

The microfiber also has a precise bipartite morphology. Each emission is a product of an additive cursor class and a multiplicative class. The three positive regions share the same additive class of eighteen conditions; together they contain three multiplicative classes of eight conditions each. The transverse region has a product \(18\times24\); the three positive regions together, another \(18\times24\); and the return, a product \(18\times8\). Decomposition by emissions has five parts, while decomposition by connected components of the bipartite graph has three. Both preserve the same \(1008\) seeds without reducing shape to a cardinality.

% Procedencia: inventario REV05, SR-015--027; ampliacion 04_tpk_emisor,
% firmas y factorizacion; X_RECIPROCIDAD/sections__generacion.tex:16--43.
% Ampliacion de la prueba de metadatos y de las tres componentes.
\subsection{Oriented metadata and decomposition of the regional microfiber}
\label{sec:microfibra-bipartita-explicita}

The additive signature preserves the diagonal class of origin and the orientation
of its shift. For an initial cursor $(i,j;d)$, with indices between
zero and eight, define
\[
 c=[i+j]_9,\qquad
 h(d)=\begin{cases}ES,&d\in\{E,S\},\\NO,&d\in\{N,O\}.
 \end{cases}
\]
In consecutive additive activations, the class $c$ increases by one
for $ES$ and decreases by one for $NO$. After the ninth activation the
direction is reversed. The sample sequence is therefore determined
by $(c,h)$, with reversal after the ninth sample. The sample
in class $c'$ is $\rho_9(c'+2)$; each decimal signature sums three consecutive
samples and reduces the total modulo ten.

\begin{proposition}[Recovery of additive metadata]
The reader $(c,h)\mapsto f_+(c,h)$ is bijective onto the eighteen
additive signatures of the emitter. Each preimage in cursor space
consists of eighteen initial conditions. The two coordinates are
recovered through the following table.
\end{proposition}

\begin{center}
\begin{tabular}{c|cc}
\toprule
$c$&$f_+(c,ES)$&$f_+(c,NO)$\\
\midrule
0&\texttt{988212}&\texttt{212988}\\
1&\texttt{212645}&\texttt{645212}\\
2&\texttt{546988}&\texttt{988546}\\
3&\texttt{889221}&\texttt{221889}\\
4&\texttt{122564}&\texttt{564122}\\
5&\texttt{465898}&\texttt{898465}\\
6&\texttt{898122}&\texttt{122898}\\
7&\texttt{221456}&\texttt{456221}\\
8&\texttt{654889}&\texttt{889654}\\
\bottomrule
\end{tabular}
\end{center}

\begin{proof}
For each entry, the class recurrence described before the
statement is applied over eighteen activations; their groups of three give the six
printed digits. For example, from $(c,h)=(4,NO)$ the first nine
classes are $4,3,2,1,0,8,7,6,5$ and the samples are
$6,5,4,3,2,1,9,8,7$. The three totals are $15,6,24$, which publish
$5,6,4$. Reversal leaves the second half $1,2,2$, yielding
\texttt{564122}. The eighteen strings in the table are distinct,
so the reader is injective and enumerates its entire image. For each
$c$ there are nine pairs $(i,j)$, one for each $i$, and for each $h$ there are
two directions. This gives $9\cdot2=18$ conditions per signature and
$18\cdot18=324$ conditions in total.
\end{proof}

Let $A_s=f_+^{-1}(s)$ and $B_e=f_\times^{-1}(e)$ be the cursor classes
of the two sheets. The decimal law recovers both signatures of each emission
$U$: if $U_j=100s_j+10t_j+u_j$, then
$s_j=\lfloor U_j/100\rfloor$ and $e_j=[t_j-s_j]_{10}$.
Consequently, the preimage of $U$ is exactly $A_s\times B_e$.
The initial conditions remain in that product; their cardinalities
are an additional reading of the same preimage.

\begin{definition}[Bipartite incidence of the microfiber]
For the oriented word $w=\texttt{010211}$, form the graph whose
vertices in one class are the additive signatures appearing over $w$ and
whose vertices in the other class are the multiplicative signatures. An
edge $(s,e)$ represents the corresponding emission and preserves as
its labeled preimage the product $A_s\times B_e$.
\end{definition}

Recovery of the signatures of the five emissions gives the complete table
of edges:
\begin{center}
\begin{tabular}{llrrl}
\toprule
$s$&$e$&$|A_s|$&$|B_e|$&stratum\\
\midrule
\texttt{564122}&\texttt{555555}&18&24&transverse\\
\texttt{898122}&\texttt{339555}&18&8&positive\\
\texttt{898122}&\texttt{393555}&18&8&positive\\
\texttt{898122}&\texttt{933555}&18&8&positive\\
\texttt{898465}&\texttt{933717}&18&8&return\\
\bottomrule
\end{tabular}
\end{center}

\begin{theorem}[Exact decomposition into three components]
The bipartite incidence has three connected components, isomorphic to
$K_{1,1}$, $K_{1,3}$ and $K_{1,1}$. Their joint preimage is the disjoint
union
\begin{align*}
 &A_{\texttt{564122}}\times B_{\texttt{555555}},\\
 &A_{\texttt{898122}}\times
  (B_{\texttt{339555}}\sqcup B_{\texttt{393555}}
     \sqcup B_{\texttt{933555}}),\\
 &A_{\texttt{898465}}\times B_{\texttt{933717}}.
\end{align*}
The respective cardinalities are $432,432,144$ and sum to $1008$.
Decomposition by emissions consists of five products and
decomposition by components consists of three sets.
\end{theorem}
\begin{proof}
The three printed additive signatures are distinct. The five multiplicative
signatures are also distinct. Only the three positive edges share
a vertex, the one associated with \texttt{898122}; they form a connected star
with three edges. The other two edges share no vertices with
it or with each other, and each forms a component. Preimages of
different signatures under a function are disjoint. Factorization of the
emitter and distributivity of Cartesian product over that disjoint
union prove the three expressions. Their cardinalities are
$18\cdot24$, $18(8+8+8)$ and $18\cdot8$. The five edges remain
individually distinguished within those components, with their multiplicity and
ordered emission.
\end{proof}

The metadata recovery map gives $(c,h)=(4,NO)$ for the transverse region,
$(6,ES)$ for each of the three positive regions and $(5,NO)$ for the return.
Thus, orientation, diagonal class, incidence and multiplicity are read from
the same signatures produced by the cursors. The distinction between five
emissions and three components preserves the complete census structure,
including the coincidence of cardinality $432$ between two distinct components.

\subsection{Finite transport and the prolongation boundary}

The endomorphisms \(L_0,L_1\) of ternary word space act before Archimedean realization. Their determination begins with the arithmetic state's transitions, not with prescribing two matrices. If \(U_t\) is the composite step of selection, transport and update, the observed block satisfies
\[
 b_j^\chi=\Pi_6(x_j^\chi),\qquad
 x_{j+1}^\chi=U_{9j+9}\cdots U_{9j+1}(x_j^\chi),
 \qquad \chi\in\{\mathrm{cl},\mathrm{pr},\mathrm{au}\}.
\]
Six transitions are formed for each regime: two consecutive transitions in each of the three regions. The pairs of source and target matrices obtained in the record are
\[
\begin{array}{c|c|c|c}
X_0&Y_0&X_1&Y_1\\\hline
010211&012222&010211&002111\\
012222&010211&002111&110221\\
201101&121221&102011&012222\\
121221&102011&012222&102011\\
121200&112202&121020&010210\\
112202&121020&010210&010200
\end{array}
\]
Each six-digit string represents a row in \(\FF_3^6\). Since \(\det X_0=1\) and \(\det X_1=-1\) in \(\FF_3\), the equations \(X_aL_a=Y_a\) uniquely determine \(L_a=X_a^{-1}Y_a\). In that row-vector chart we obtain
\[
 L_0=\begin{pmatrix}
 2&2&2&1&2&1\\2&1&2&2&1&1\\1&0&0&1&0&2\\
 0&0&1&1&1&0\\2&2&2&0&2&1\\2&1&2&1&0&0
 \end{pmatrix},\quad
 L_1=\begin{pmatrix}
 2&1&2&0&2&2\\1&2&1&1&2&0\\1&2&1&1&1&2\\
 0&1&0&0&2&1\\2&0&2&1&0&1\\0&2&2&2&1&1
 \end{pmatrix}\quad(\bmod3).
\]
The calendar \((L_0,L_0,L_1,L_1)\) produces four new six-component blocks. Their concatenations form
\[
 w_6\longrightarrow w_{12}\longrightarrow w_{18}
 \longrightarrow w_{24}\longrightarrow w_{30}.
\]
The subscripts of \(w\) indicate accumulated length. The boundary \(R_{36}\) preserves the terminal structure and opens the extension relation; it is not renamed as a final periodic block that could be repeated indefinitely.

The identity \(L=X^{-1}Y\) proves uniqueness once the transitions are fixed; it does not replace production of those transitions by \(U_t\). The provenance of their register and examination of the sixteen binary calendars are documented in the technical supplement. The indicated calendar is the only one of those sixteen that jointly reproduces the three sequences of five blocks.

\subsection{Five closure regions and angular compensation}

The closure region contains one transverse component and four companions. The symmetric reader acts on the space of its five positions through
\[
 P_5=\frac15\mathbf1\mathbf1^{\mathsf T},\qquad
 \mathbf1=(1,1,1,1,1)^{\mathsf T}.
\]
Invariance \(P_5\lambda=\lambda\) and conservation \(\mathbf1^{\mathsf T}\lambda=1\) force \(\lambda=\mathbf1/5\). This normalization distributes unity among reader positions; it does not replace seed multiplicities \(432,144,144,144,144\) with uniform frequencies. The reader's elliptic chart realizes each coordinate through \(1+\lambda_jJ\); its slope is \(\lambda_j\), so the primitive slope is \(1/5\). The rule passing from a normalized coordinate to a slope is thus made explicit, instead of identifying the two notions without a map. The reader composes the four companion positions by the law
\begin{equation}
 a\oplus b=\frac{a+b}{1-ab}.
 \label{x_reciprocidad:eq:suma-pendientes}
\end{equation}
This law is obtained by multiplying \((1+aJ)(1+bJ)\) in regime \(J^2=-I\) and dividing by its scalar component. Composition of slopes is therefore a realization of TRIT's quadratic structure, not an ordinary sum of four fractions.

\begin{proposition}[Closure compensator]
Composition of four slopes \(1/5\) is \(120/119\). The condition of return to slope one determines a unique positive compensator of the form \(1/q\), with \(q>1\), and gives \(q=239\).
\end{proposition}
\begin{proof}
The first doubling gives \((1/5)\oplus(1/5)=5/12\); the second, \((5/12)\oplus(5/12)=120/119\). The condition
\[
 \frac{120/119-1/q}{1+(120/119)/q}=1
\]
is equivalent to \((120-119)q=120+119\), so \(q=239\).
\end{proof}

The resulting closure character has realization
\begin{equation}
 \Pi_{\rm cl}=4\left(4\arctan\frac15-\arctan\frac1{239}\right).
 \label{x_reciprocidad:eq:lector-pi}
\end{equation}
Machin's angular identity recognizes \(\Pi_{\rm cl}=\pi\). The construction direction is precise: the slope and number of compositions are declared by the pentafiber reader; the compensation equation forces \(239\). Once this character is produced, its evaluation by alternating series does not select the seed orbit anew.

For \(0<x<1\), the alternating sums
\[
 A_n(x)=\sum_{j=0}^n\frac{(-1)^jx^{2j+1}}{2j+1}
\]
enclose \(\arctan x\) between two consecutive truncations, with difference \(x^{2n+3}/(2n+3)\). Applied to the two slopes of \eqref{x_reciprocidad:eq:lector-pi}, they provide rational endpoints with explicit error. The numerical value is thus calculated as a consequence of the exact character, without introducing a target prefix.

\subsection{Propagation reader and normalization of unity}

The propagation region is realized through a formal family \(P_t\in\mathcal A[[t]]\), where \(\mathcal A\) is a unital commutative algebra over \(\QQ\). Its composition satisfies \(P_{t+s}=P_tP_s\), with unit \(P_0=1\), and oriented elementary transport fixes the scalar generator \(P'_0=+1\). The condition fixes its sign, not just its magnitude. Writing
\[
 P_t=\sum_{n\ge0}a_nt^n,
\]
normalization determines \(a_0=a_1=1\). Comparing the coefficient of \(s^nt\) in \(P_{s+t}=P_sP_t\) gives \((n+1)a_{n+1}=a_na_1\); equivalently, \(P'_t=P_t\). Thus,
\begin{equation}
 (n+1)a_{n+1}=a_n,\qquad a_n=\frac1{n!}.
 \label{x_reciprocidad:eq:propagacion}
\end{equation}
The character at one parameter unit is \(\Pi_{\rm pr}=P_1\), subsequently recognized as \(e\).

Normalization of the generator matters. Without that datum the semigroup law would admit \(P_t=e^{ct}\) for different \(c\); hence the condition fixing the parameter unit in the reader is not omitted. This condition precedes decimal comparison and belongs to exposition of the operational rule, not to selection of an experimental value.

\begin{proposition}[Rational propagation bounds]
For \(n\ge1\), let \(E_n=\sum_{j=0}^n1/j!\). Then
\[
 E_n<\Pi_{\rm pr}<E_n+\frac{n+2}{(n+1)(n+1)!}.
\]
\end{proposition}
\begin{proof}
The first omitted term is \(1/(n+1)!\). Each subsequent ratio is at most \(1/(n+2)\). The geometric sum majorizing the tail is
\[
 \frac1{(n+1)!}\frac1{1-1/(n+2)}
 =\frac{n+2}{(n+1)(n+1)!}.
\]
The strict inequality follows from the decrease of successive ratios.
\end{proof}

\subsection{Modal incidence and generation of self-scaling}

The tritic selector produces the cycle \((+1,-1,0)\). In modes \(\Sigma,\Pi\), the rule is \(+1:m\mapsto\Sigma\), \(-1:m\mapsto\Pi\), \(0:m\mapsto m\). With cyclic closure, the observed mode is \((\Sigma,\Pi,\Pi)\). Its edges are \(\Sigma\to\Pi\), \(\Pi\to\Pi\) and \(\Pi\to\Sigma\). Its incidence, in that mode order, is therefore
\begin{equation}
 F_{\rm av}=\begin{pmatrix}0&1\\1&1\end{pmatrix}.
 \label{x_reciprocidad:eq:fav}
\end{equation}
The four entries are obtained from the calendar, not from an equation that previously received the value of \(\varphi\). Repetition of the calendar gives counts \(3F_{\rm av},9F_{\rm av},36F_{\rm av}\) at lengths nine, twenty-seven and one hundred and eight. These counts are not powers of the matrix.

\begin{theorem}[Positive self-scaling character]
The incidence \eqref{x_reciprocidad:eq:fav} has a unique strictly positive eigenray. Normalizing it as \((1,x)^{\mathsf T}\), its coordinate \(x\) satisfies \(x^2-x-1=0\), \(x>0\), and is \(\varphi=(1+\sqrt5)/2\).
\end{theorem}
\begin{proof}
The eigenvalue equation gives \(x=\lambda\) and \(1+x=\lambda x\). Eliminating \(\lambda\) produces the polynomial. Only one of its two roots is positive. Since \(F_{\rm av}^2\) has all entries positive, the positive ray is unique and simple. The radical expression recognizes the resulting coordinate; it does not participate in production of \(F_{\rm av}\).
\end{proof}

With \(F_0=0,F_1=1,F_{n+1}=F_n+F_{n-1}\), the incidence induces, for \(n\geq1\),
\[
 F_{\rm av}^n=\begin{pmatrix}F_{n-1}&F_n\\F_n&F_{n+1}\end{pmatrix}.
\]
The consecutive ratios \(F_{n+1}/F_n\) alternate around \(\varphi\). The determinant identity
\[
 F_{n+1}^2-F_nF_{n+2}=(-1)^n
\]
gives an exact width \(1/(F_nF_{n+1})\) between two consecutive ratios. Their growth makes this rational bracket arbitrarily small.

The measure on all compatible histories does not coincide with the mode frequency of the three-event calendar. The chronological frequency is \((1/3,2/3)\). The measure of maximal entropy of the symbolic envelope defined by \(F_{\rm av}\) is constructed from its eigenvectors and has weight ratio \(\varphi^{-2}\). The latter enters the areal--volumetric reading; it is not obtained merely by replacing one counting frequency with another.


% BEGIN OWNER /Users/ruben/Documents/New project/output/EXTREMA_MEDIA_RAZON_RECIPROCIDAD_ES_EN_20260918/ARTICULO/ES/source/sections/retorno_areal_volumetrico.tex
\subsubsection{Trajectory space and invariant measure between sheets}
\label{x_reciprocidad:sec:medida-retorno-hojas}

The toroidal topology of APP fixes the support of cardinal translations
and their winding classes. The correspondence between sheets comes from
another operation on that same support: temporal selection of the
additive or multiplicative evaluation. Its geometric typing is
\[
 \eta_{\rm av}(\Sigma)=\mathscr H_{\rm ar},\qquad
 \eta_{\rm av}(\Pi)=\mathscr H_{\rm vol}.
\]
The first sheet corresponds to the additive boundary reading; the second,
to the multiplicative interior reading. This map transports the mode
labels and preserves the incidence \eqref{x_reciprocidad:eq:fav}. The square of the self-scale
ratio is determined through the trajectories of that incidence.

\Needspace{8\baselineskip}
\begin{definition}[Memory-one symbolic envelope]
Let
\[
 X_{\rm av}=\{(s_k)_{k\in\mathbb Z}\in
 \{\mathscr H_{\rm ar},\mathscr H_{\rm vol}\}^{\mathbb Z}:
 (F_{\rm av})_{s_k,s_{k+1}}=1\text{ for every }k\}.
\]
The index shift acts on this space. It is the smallest memory-one symbolic
system containing the modal sequences and preserving
their consecutive pairs: it admits both cross transitions and
volumetric self-retention. Phase, cardinal orientation and trajectory
registers remain in the TPK state from which this
envelope is obtained.
\end{definition}

\begin{proposition}[Weighting of trajectories between sheets]
In areal--volumetric order, the invariant measure of maximal entropy
of \(X_{\rm av}\) has transition matrix and stationary distribution
\begin{equation}
 P_{\rm av}=\begin{pmatrix}0&1\\\varphi^{-2}&\varphi^{-1}\end{pmatrix},
 \qquad
 (\mu_{\rm ar},\mu_{\rm vol})
 =\frac{(1,\varphi^2)}{1+\varphi^2}.
 \label{x_reciprocidad:eq:medida-hojas-exacta}
\end{equation}
In particular, \(\mu_{\rm ar}/\mu_{\rm vol}=\varphi^{-2}\).
\end{proposition}
\begin{proof}
The already constructed incidence has positive eigenvector
\(r=(1,\varphi)^{\mathsf T}\) and eigenvalue \(\varphi\). The
normalization
\[
 (P_{\rm av})_{ij}
 =\frac{(F_{\rm av})_{ij}r_j}{\varphi r_i}
\]
gives the indicated matrix. Its rows sum to one by
\(\varphi^{-2}+\varphi^{-1}=1\). Symmetry of \(F_{\rm av}\)
implies that weights proportional to \(r_i^2\) satisfy detailed
balance; therefore, \(\mu P_{\rm av}=\mu\). Irreducibility and
volumetric self-retention determine a unique stationary distribution.

To verify maximality, every allowed transition satisfies
\(\log(P_{\rm av})_{ij}=\log r_j-\log r_i-\log\varphi\).
Under any invariant measure, the endpoint terms cancel
on average. If \(\nu_i\) are its vertex weights and
\(q_i(j)=\nu(s_1=j\mid s_0=i)\), for \(\nu_i>0\), then
\[
 h_\nu\leq H_\nu(s_1\mid s_0)
 =\log\varphi-\sum_{i:\nu_i>0}\nu_iD(q_i\Vert P_i)
 \leq\log\varphi,
\]
where \(D(q\Vert p)=\sum_jq_j\log(q_j/p_j)\) is the relative
entropy on allowed pairs and \(0\log0=0\) is adopted.
The Markov measure of \eqref{x_reciprocidad:eq:medida-hojas-exacta} attains the bound.
Equality in the first inequality requires the Markov property;
equality in the second requires \(q_i=P_i\). Stationarity and
irreducibility then fix \(\mu\), establishing
uniqueness. This is the Parry measure of the defined envelope.
\end{proof}

The frequency \((1/3,2/3)\) counts the instants of the primitive calendar;
\(\mu\) weights the trajectories of the envelope. Its quadratic ratio
expresses the product of incoming and outgoing incidence amplitudes.
TPK memory and phase restrictions retain their own
domain; the envelope's entropy corresponds to the symbolic system
specified in this definition.

\subsubsection{Closure operator and limit of marked returns}
\label{x_reciprocidad:sec:operador-retorno-marcado}

The closure coordinate \(\pi\), generated by the regional reader,
enters after construction of modal incidence. Let
\(e_{\rm ar}=(1,0)^{\mathsf T}\),
\(e_{\rm vol}=(0,1)^{\mathsf T}\) and
\(P_i=e_ie_i^{\mathsf T}\) be the sheet projectors. The reader's conditions
are preservation of both sheets, unit volumetric normalization
and areal mark equal to the closure coordinate. These conditions
determine the positive operator
\begin{equation}
 D_\pi=\pi P_{\rm ar}+P_{\rm vol}
       =\begin{pmatrix}\pi&0\\0&1\end{pmatrix}.
 \label{x_reciprocidad:eq:operador-marca-pi}
\end{equation}
Diagonality expresses preservation of sheets and the two entries
express their normalizations; uniqueness of the operator
for the fixed reader is thus explicit.

\Needspace{8\baselineskip}
\begin{proposition}[Ratio of boundary returns]
For \(n\geq1\), the marked evaluation
\begin{equation}
 \Theta_n=
 \frac{\langle e_{\rm ar},D_\pi F_{\rm av}^{\,n}e_{\rm ar}\rangle}
 {\langle e_{\rm vol},F_{\rm av}^{\,n}e_{\rm vol}\rangle}
 =\pi\frac{F_{n-1}}{F_{n+1}}
 \label{x_reciprocidad:eq:retorno-marcado-finito}
\end{equation}
satisfies
\begin{equation}
 \lim_{n\to\infty}\Theta_n
 =\pi\frac{\mu_{\rm ar}}{\mu_{\rm vol}}
 =\frac{\pi}{\varphi^2}.
 \label{x_reciprocidad:eq:retorno-marcado-limite}
\end{equation}
\end{proposition}
\begin{proof}
The diagonal entries of \(F_{\rm av}^{\,n}\) count paths
of length \(n\) returning to their initial sheet. The already proved power
formula provides \(F_{n-1}\) and \(F_{n+1}\). The operator
\(D_\pi\) multiplies the first count by \(\pi\). Finally,
\(F_{n-1}/F_{n+1}\) is the product of two consecutive ratios
converging to \(\varphi^{-1}\); its limit is
\(\varphi^{-2}\). The identity with the weight ratio follows from
\eqref{x_reciprocidad:eq:medida-hojas-exacta}.
\end{proof}

The closure operator appears only once, in the terminal evaluation
of the return. The transfer \(F_{\rm av}\) remains determined
by the calendar. This composes information from two readers of the
same state: modal incidence determines the quadratic weighting
and the closure region contributes its coordinate to the boundary.

The composition also preserves refinement. For positive cylinders
\(I=[a,b]\) and \(J=[c,d]\), the image under the reader
\(\mathcal R(x,y)=x/y^2\) is the interval
\[
 \mathcal Q(I,J)=\left[\frac{a}{d^2},\frac{b}{c^2}\right].
\]
Monotonicity in each argument proves that the nested cylinders of
closure and self-scale produce nested intervals of the ratio;
positivity of their limits and contraction of their diameters determine
\(\pi/\varphi^2\). The scalar expression thus retains a chain
of operations compatible with depth, subsequently applied in
its two orientations.

% END OWNER


After the characters are generated, the two orientations of the reader
\[
 R(x,y)=\frac{x}{y^2}
\]
provide
\begin{equation}
 \rho_A=\frac\pi{\varphi^2},\qquad
 \rho_E=\frac\varphi{\pi^2},\qquad
 \varphi^3\rho_A^2\rho_E=1.
 \label{x_reciprocidad:eq:areal-electron}
\end{equation}
The second expression enters the exponent of the electronic construction in Section~\ref{x_reciprocidad:el-centro:registros}. The first is also obtained as the limit of the marked reading \(\pi F_{n-1}/F_{n+1}\). Exchanging arguments relates the two, but does not turn a reparametrization of the electronic equation into a second independent derivation of all its coefficients.


% BEGIN OWNER /Users/ruben/Documents/New project/output/EXTREMA_MEDIA_RAZON_RECIPROCIDAD_ES_EN_20260918/ARTICULO/ES/source/sections/euler_trit_articulacion.tex
% Articulación posterior a la generación de pi, phi y e.
% Contenido acumulativo; no reemplaza el desarrollo previo del TRIT.
\subsection{Exponential realizations of the generated coordinates}
\label{x_reciprocidad:euler-trit-articulacion}

The already determined regional coordinates allow the specific
realizations of TRIT's quadratic family to be recognized. The order is
substantive: construction of the regions precedes evaluation of an
exponential at parameters \(\pi\) or \(2\log\varphi\). These expressions
identify the function of the coordinates in operators constructed from
APP and TPK; they do not enter into selection of the prefixes producing them.

\subsubsection{Ternary cycle, nilpotent transport and sheet incidence}

Multiplication \(a(r)=4r\) in \(\mathbb Z/9\mathbb Z\) has order
\(3\). On the units it determines the orbits \(\{1,4,7\}\) and
\(\{2,8,5\}\). In the augmentation module of either, formed by
integer combinations whose coefficients sum to zero, the basis
\((e_r-e_{4r},e_{4r}-e_{7r})\) gives
\[
 C=\begin{pmatrix}0&-1\\1&-1\end{pmatrix},
 \qquad C^2+C+I=0.
\]
The restricted form has Gram matrix
\(\bigl(\begin{smallmatrix}2&-1\\-1&2\end{smallmatrix}\bigr)\).
In residue evaluations modulo \(9\), the action distinguishes the sheets:
\(S(ax,ay)=aS(x,y)\), whereas \(P(ax,ay)=a^{-1}P(x,y)\).
The cyclic realization therefore preserves a contragredient orientation
between sum and product. The quotients of the integer lift are preserved
in the corresponding arithmetic register.

Elementary parabolic transport is represented by
\[
 N=\begin{pmatrix}0&1\\0&0\end{pmatrix}.
\]
The areal--volumetric incidence already constructed in \eqref{x_reciprocidad:eq:fav} provides
\[
 F_{\rm av}=\begin{pmatrix}0&1\\1&1\end{pmatrix},
 \qquad A=F_{\rm av}^2=\begin{pmatrix}1&1\\1&2\end{pmatrix}.
\]
The three operators act on declared real spaces: the augmentation
module tensored with \(\mathbb R\), the nilpotent-transport plane and
the modal-incidence plane. Sharing matrix dimension does not identify
these spaces or their coordinates.

\begin{proposition}[Concrete generators of the three types]
\label{x_reciprocidad:euler-trit-generadores}
The operators
\[
 J_{\mathrm e}=\frac{2C+I}{\sqrt3},\qquad
 J_{\mathrm p}=N,\qquad
 J_{\mathrm h}=\frac{2A-3I}{\sqrt5}
\]
satisfy, respectively,
\[
 J_{\mathrm e}^{\,2}=-I,\qquad
 J_{\mathrm p}^{\,2}=0,\qquad
 J_{\mathrm h}^{\,2}=I.
\]
\end{proposition}
\begin{proof}
The relation in \(C\) implies \((2C+I)^2=-3I\).
Direct multiplication gives \(N^2=0\).
Finally, \(A\) has trace \(3\) and determinant \(1\), so
\(A^2-3A+I=0\) and \((2A-3I)^2=5I\).
\end{proof}

\begin{proposition}[Three evaluations of the exponential identity]
\label{x_reciprocidad:euler-trit-evaluaciones}
In the preceding realizations,
\begin{equation}
 \boxed{
 \exp(\pi J_{\mathrm e})+I=0,\qquad
 \exp(J_{\mathrm p})=I+J_{\mathrm p},\qquad
 \exp(2\log\varphi\,J_{\mathrm h})=F_{\rm av}^{\,2}.}
 \label{x_reciprocidad:euler-trit-tres-evaluaciones}
\end{equation}
\end{proposition}
\begin{proof}
The elliptic specialization of the tritic exponential gives
\(\cos\pi\,I+\sin\pi\,J_{\mathrm e}=-I\).
Nilpotency removes exactly the terms of degree at least \(2\)
in the second expression. For the third, \(\varphi^2=\varphi+1\) implies
\[
 \cosh(2\log\varphi)=\frac32,\qquad
 \sinh(2\log\varphi)=\frac{\sqrt5}{2}.
\]
Therefore, the exponential equals
\(\frac32I+\frac{\sqrt5}{2}J_{\mathrm h}=A\).
\end{proof}

The first equality represents the elliptic half-turn; the full turn
satisfies \(\exp(2\pi J_{\mathrm e})=I\). The second retains a nonzero linear
term: the parabolic regime expresses transport, not absence of
evolution. The third realizes a unimodular hyperbolic advance. Its
eigenvalues are \(\varphi^2\) and \(\varphi^{-2}\); one direction expands and
the other contracts while preserving the determinant. The three evaluations
use distinct parameters of a common exponential family.

In that family, \(e\) enters through the scalar evaluation
\(\exp(I)=eI\) and the law \(\exp(sI)\exp(tI)=\exp((s+t)I)\).
Its role is not restricted to the parabolic type. The coordinate \(e\), previously
obtained from its region, is recognized here as the unit evaluation
of the exponential; this identity does not replace its coinductive genealogy.

\subsubsection{Polynomial resolution of the regimes}

The three realizations admit a compact presentation on the direct sum
of their spaces. Define
\[
 \mathbb J=J_{\mathrm e}\oplus J_{\mathrm p}\oplus J_{\mathrm h},
 \qquad Q=\mathbb J^2.
\]
The preceding relations imply
\[
 \mathbb J^6=\mathbb J^2,\qquad Q^3=Q.
\]
Use of \(Q\) keeps this operator square separate from the dodecaphasic
register \(K\).

\begin{proposition}[Polynomial regime projectors]
\label{x_reciprocidad:euler-trit-proyectores}
The maps
\[
 p_{\mathrm e}=\frac{Q^2-Q}{2},\qquad
 p_{\mathrm p}=I-Q^2,\qquad
 p_{\mathrm h}=\frac{Q^2+Q}{2}
\]
are mutually orthogonal idempotents, sum to \(I\) and recover the three
summands of the representation.
\end{proposition}
\begin{proof}
The three polynomials take values \((1,0,0)\), \((0,1,0)\) and
\((0,0,1)\) at the points \(-1,0,+1\), respectively. The differences
between their squares and themselves, as well as their cross products, are
multiples of \(X^3-X\). Substituting \(Q\), which annihilates that polynomial, proves
the identities. On the three blocks, \(Q\) acts as \(-I,0,I\).
\end{proof}

We thus obtain
\begin{align}
 \exp(t\mathbb J)
 ={}&p_{\mathrm e}(\cos t\,I+\sin t\,\mathbb J)
   +p_{\mathrm p}(I+t\mathbb J)\notag\\
   &+p_{\mathrm h}(\cosh t\,I+\sinh t\,\mathbb J).
 \label{x_reciprocidad:euler-trit-exponencial-total}
\end{align}
Each summand preserves its quadratic relation and conjugation reverses
its evolution. The total representation combines the three regimes, but does not
by itself define a transition operator between them. Such a transition
requires the TPK transport maps with their domains, orientation and
memory. Exponential articulation allows their types to be recognized without
replacing that dynamics with identification of the fibers.

% PROCEDENCIA FOCAL (no impresa)
% Resultado recuperado: tres generadores y tres evaluaciones exponenciales.
% Propietario completo:
% BASE/manuscrito/sections/hmt/09b_algebras_cuadraticas.tex:289-516.
% Antecedente causal de C y acción contragrediente:
% BASE/manuscrito/sections/hmt/03d_esqueleto_ciclotomico_app.tex:13-47,145-246.
% Los tres tipos y la orientación están en:
% BASE/manuscrito/sections/hmt/02_trit_geometrias.tex:83-196.
% Resolución polinómica recuperada:
% output/INVESTIGACION_HOLONOMIA_NONADICA_CRISTAL_APERIODICO_20260903/
% ALGEBRA_HOLONOMICA_UNIVERSAL_Y_EULER_TRITICO.md, secciones 8-11 y 22.
% Los cuatro propietarios se leyeron completos. BASE designa:
% /Users/ruben/Documents/HMT2/
% EL_CIERRE_HOLOGRAFICO_DEL_INFINITO_SUCESOR_102_CAPITULOS_2026-08-28_EN_TRABAJO/
% No se incorpora la identidad P9(alpha)=delta del documento especializado:
% esa identificación no pertenece a este bloque.
% Tampoco se identifica una dirección nilpotente con la vacancia electrónica,
% ni se presume una representación global del transporte por la suma directa.
% COMPROBACIONES: Python estándar, enteros y Fraction.
% C^2+C+I=0; (2C+I)^2=-3I; N^2=0; (2Fav^2-3I)^2=5I.
% 81 pares de residuos verifican las dos acciones contragredientes.
% Los 3 proyectores se verificaron en Q=-1,0,+1; prueba general en texto.
% 6 labels únicos y entornos equilibrados. Sin compilación.

% END OWNER


\subsection{Compatible cylinders and arbitrary depth}

The realization of a ternary block \(b=(b_1,\ldots,b_6)\) is the integer
\[
 \nu(b)=\sum_{j=1}^6b_j3^{6-j}\in\{0,\ldots,728\}.
\]
A sequence of blocks defines integer prefixes by \(N_{n+1}=729N_n+\nu(b_{n+1})\), starting at \(N_0=m\), where \(m\) is the integer part determined by the reader. The preceding brackets place the closure, propagation and self-scale characters in \((3,4)\), \((2,3)\) and \((1,2)\), respectively; their initial values are therefore \(m=3,2,1\). Subsequent blocks refine the fractional part. The corresponding cylinder is
\begin{equation}
 I_n=\left[\frac{N_n}{729^n},\frac{N_n+1}{729^n}\right],
 \qquad |I_n|=729^{-n}.
 \label{x_reciprocidad:eq:cilindro}
\end{equation}
The half-open convention is used to select a representative when a value falls on a boundary. Passage to the limit uses the closures of these intervals and identifies the two equivalent boundary expansions. This precision prevents assuming that every intersection of nested half-open intervals is automatically nonempty.

\begin{theorem}[Determination of the scalar realization]
A compatible block history determines a unique limit of the left endpoints of \eqref{x_reciprocidad:eq:cilindro}. If an exact character has rational brackets of arbitrarily small width and is separated from the boundary of each decimal cylinder examined, each decimal prefix of finite length is determined after a finite calculation. The prefixes thus obtained are mutually compatible.
\end{theorem}
\begin{proof}
Extension by a block preserves the preceding prefix by Euclidean division. The closed cylinders are nested and their diameter tends to zero; their left endpoints form a Cauchy sequence and have a unique limit. For a fixed decimal length, a value separated from the partition boundaries has positive distance from them. A sufficiently narrow bracket lies within a unique cell and determines its prefix. Uniqueness of the cell and nesting of the partitions prove compatibility under truncation. At points with a terminating expansion, the exact boundary decision and its representation convention must be added; that decision does not follow from a numerical width alone.
\end{proof}

\begin{corollary}[Finite determination and compatibility of regional prefixes]
\label{x_reciprocidad:gen:prefijos-regionales}
For each of the realizations \(\pi,\varphi,e\) and every finite decimal
length, the constructed brackets determine the corresponding prefix
through a finite calculation. Prefixes obtained at different depths
are compatible under truncation.
\end{corollary}
\begin{proof}
The closure, propagation and self-scale characters have the preceding
brackets. After their recognition, irrationality of \(\pi,e,\varphi\)
guarantees their separation from the rational boundaries of each finite decimal
partition. The preceding theorem is applied to each character.
\end{proof}
Arbitrary depth expresses the quantifier of the corollary: for every
finite horizon there is a calculation determining its prefix. The mathematical
object is the complete compatible family; each run obtains a
finite projection of it.

\subsection{Mirror involution and synchronization of the regions}

The nonadic connection jointly transports the three characters. In the ninth-phase chart, the self-scale axis and ternary sum of closure and propagation remain fixed under mirror exchange. Helical inversion changes the transport direction; the two-channel mirror changes their relative orientation. They are different involutions, although both act on the same trajectory system.

Census coincidences illustrate this coordination:
\begin{center}
\begin{tabular}{@{}c rrr l@{}}\toprule
Phase&\(N_\pi\)&\(N_e\)&\(N_\varphi\)&Equality of counts\\\midrule
4&207&207&122&\(N_\pi=N_e\)\\
5&39&151&151&\(N_e=N_\varphi\)\\
6&110&110&69&\(N_\pi=N_e\)\\
7&81&29&81&\(N_\varphi=N_\pi\)\\
8&59&59&59&triple synchronization\\
9&43&43&43&triple synchronization\\\bottomrule
\end{tabular}
\end{center}
These integers count compatible coordinates of the regions. They do not assert equality between \(\pi\), \(e\) and \(\varphi\) as real numbers. If
\[
 \partial(a,b,c)=(a-b,b-c,c-a),
\]
its image belongs to the lattice \(A_2=\{(u,v,w)\in\ZZ^3:u+v+w=0\}\). In phases eight and nine the relative component vanishes, but the diagonal component changes from \((59,59,59)\) to \((43,43,43)\). Relative synchronization coexists with a shift of \(-16\) per channel. The closure leading to \(\alpha\) must preserve both pieces of information: relative coincidence and evolution of memory.

\subsection{Sturmian helical suspension}

The relation between bases nine and ten of the completion defines the dimensionless parameter
\begin{equation}
 \delta=\frac{\log(10/9)}{\log10}.
 \label{x_reciprocidad:eq:delta-sturm}
\end{equation}
It is not obtained from a target physical constant. It is the relative logarithmic scale of two bases already present in the construction. Its use as a symbolic reading of refinement has an exact consequence.

\begin{lemma}[Irrationality of the slope]
The number \(\delta\) is irrational.
\end{lemma}
\begin{proof}
If \(\delta=p/q\), \(q>0\), then \((10/9)^q=10^p\) and \(9^q=10^{q-p}\). Factorization of the left-hand side contains the prime three and that of the right-hand side only the primes two and five. The equality is impossible.
\end{proof}

The lower mechanical word, with intercept \(\theta\), is
\[
 z_n(\theta)=\lfloor(n+1)\delta+\theta\rfloor-\lfloor n\delta+\theta\rfloor\in\{0,1\}.
\]
It is obtained by observing the rotation \(\theta\mapsto\theta+\delta\pmod1\) through the interval \([1-\delta,1)\). The event sequence has density \(\delta\), and the sum over any window of length \(L\) equals \(\lfloor L\delta\rfloor\) or \(\lceil L\delta\rceil\). Since \(21<1/\delta<22\), event gaps are twenty-one or twenty-two. The word indicating which intervals are short is another mechanical word of slope \(\sigma=22-1/\delta\); its events are separated by six or seven primary intervals. The second scale is induced by the first, not a second free parameter.


% BEGIN OWNER /Users/ruben/Documents/New project/output/EXTREMA_MEDIA_RAZON_RECIPROCIDAD_ES_EN_20260918/ARTICULO/ES/source/sections/vacancias_capacidad.tex
% Desarrollo reunido de DESARROLLO_MATEMATICO, apartados 38--40 y 50.
\subsubsection{Capacity index, vacancies, and induction of the tritic regime}
\label{x_reciprocidad:vac-capacidad-trit}

The vacancy sequence is linked to refinement through an
increment identity. To show it, fix the phase origin
\(\theta=0\). Let
\[
 \rho=\log_{10}9=1-\delta,\qquad
 \mathcal C_t=\lfloor(t+1)\rho\rfloor\quad(t\geq0).
\]
The index \(\mathcal C_t\) expresses the decimal capacity of refinement
in the normalization comparing a ternary block of \(6\) positions,
with \(729\) possibilities, and a decimal block of \(3\) positions, with
\(1000\) possibilities. Indeed,
\(\log_{1000}729=\log_{10}9=\rho\). The letter \(K\) is reserved for
the dodecaphasic register to be constructed later.

\begin{proposition}[Complementarity of capacity and vacancy]
\label{x_reciprocidad:prop:capacidad-vacancia}
For \(t\geq1\),
\begin{equation}
 \mathcal C_t-\mathcal C_{t-1}=1-z_t(0).
 \label{x_reciprocidad:eq:capacidad-vacancia}
\end{equation}
Consequently, if \(0\leq a<b\),
\[
 \mathcal C_b-\mathcal C_a+
 \sum_{t=a+1}^{b}z_t(0)=b-a.
\]
\end{proposition}
\begin{proof}
Irrationality of \(\delta\) gives
\[
 \mathcal C_t
 =t+1-\lceil(t+1)\delta\rceil
 =t-\lfloor(t+1)\delta\rfloor.
\]
Subtracting two consecutive terms produces
\eqref{x_reciprocidad:eq:capacidad-vacancia}. The sum over the interval telescopes.
\end{proof}

Each transition increases the block counter by one unit. If
\(z_t=0\), capacity also increases; if \(z_t=1\), capacity
remains unchanged while a transition is incorporated into the history. Vacancy
is here an event of capacity retention, with a determined position and
antecedent. The preceding balance exactly preserves length
between acquired capacity and recorded vacancies. The observed
phase \(g_t^{\rm cap}=[\mathcal C_t]_9\) satisfies
\[
 g_t^{\rm cap}-g_{t-1}^{\rm cap}\equiv1-z_t\pmod9.
\]
This capacity phase and the chronological phase \([t]_9\) are two
different observations. Local capacity retention and a complete
holonomy turn have distinct mechanisms, although both
allow phase observation to be preserved while the state changes.

The relation to TRIT can be followed in the vacancy positions themselves.
For \(j\geq1\), define
\[
 p_j=\left\lfloor\frac j\delta\right\rfloor,\qquad
 v_j=\mathcal C_{p_j},\qquad h_j=p_{j+1}-p_j.
\]
The inequality \(p_j\delta<j<(p_j+1)\delta\) identifies \(p_j\)
as the position of event \(j\). Since \(0<\delta<1\), it also gives
\(\lfloor(p_j+1)\delta\rfloor=j\), and therefore
\[
 v_j=p_j-j,\qquad
 v_{j+1}-v_j=h_j-1.
\]
If \(s_j=22-h_j\), then \(s_j\in\{0,1\}\) and
\begin{equation}
 [v_{j+1}]_3=[v_j-s_j]_3.
 \label{x_reciprocidad:eq:vacancia-regimen}
\end{equation}
An interval of length \(22\) preserves the class; one of length
\(21\) shifts it one unit in the negative direction. The canonical selector
is applied to the positive phase \(r_j=1+[v_j]_9\):
\[
 \tau_j=M_{\rm ph}(r_j).
\]
A tritic regime is thus determined by the capacity index
retained at each event, in addition to binary encoding of the vacancy.

The short-interval positions form the induced mechanical encoding
of slope \(22-1/\delta\). Their gaps \(6\) and \(7\)
determine the lengths of the plateaus of class \([v_j]_3\), once
the origin and endpoint convention are fixed. The first classes
are \(2\) for \(6\) events, \(1\) for \(7\) and \(0\) for
\(7\); the first segment of \(6\) is the initial restriction of a plateau.
Their signatures run through \(0,-1,+1\). Equations
\eqref{x_reciprocidad:eq:capacidad-vacancia} and \eqref{x_reciprocidad:eq:vacancia-regimen} make explicit the
link between refinement, vacancies and regime selection. The quadratic
TRIT reading then realizes these types through their respective
generators; the temporal event and its geometric representation
thus retain an identifiable dependence.

Inverting the monotone index provides another useful observation:
\[
 \mathcal T(k)=\min\{t:\mathcal C_t=k\}
   =\left\lceil\frac k\rho\right\rceil-1\qquad(k\geq1).
\]
If \(\sigma=\rho^{-1}-1\), the number of blocks needed to gain
\(m\) units of capacity is
\[
 \mathcal T(k+m)-\mathcal T(k)
 =m+\lceil(k+m)\sigma\rceil-\lceil k\sigma\rceil
 \in\{m+\lfloor m\sigma\rfloor,m+\lceil m\sigma\rceil\}.
\]
In particular, \(9\) units of capacity require \(9\) or \(10\)
blocks. The period of a finite phase remains exact, while the
return time measured by the other index admits two values. This
distinction is part of refinement dynamics, not an a posteriori correction
of its numerical values.

% END OWNER


\begin{theorem}[Complexity of phase readings]
For \(m\ge1\), the mechanical word has factor complexity \(p(m)=m+1\). Retaining phase \(n\bmod9\) gives complexity \(p_9(m)=9(m+1)\); retaining only its ternary class gives \(p_3(m)=3(m+1)\). All three topological entropies are zero.
\end{theorem}
\begin{proof}
For factors of length \(m\), the predecessors of the two endpoints of the coding interval form \(m+1\) distinct points on the circle: coincidence between shifted endpoints removes duplications and irrationality prevents additional coincidences. These points determine \(m+1\) intervals with constant and distinct itineraries. Every irrational orbit visits each interval, whence \(p(m)=m+1\).

With a phase modulo nine fixed, positions of that phase follow a rotation of step \(9\delta\), also irrational. By density, all factors appear in each phase, and the first coordinate distinguishes their nine copies. The same argument with step \(3\delta\) gives three copies after ternary reduction. Finally, \(m^{-1}\log(c(m+1))\to0\) for \(c=1,3,9\).
\end{proof}

The phase product \(C_9\times\mathbb T\), translated by \((1,\delta)\), has frequency characters
\[
 \frac a9+k\delta\pmod1,\qquad a\in\ZZ/9\ZZ,\ k\in\ZZ.
\]
The spectral module is not \(\log10\): it is the additive module generated by \(1/9\) and \(\delta\), modulo the integers. The factor \(\log10\) belongs to normalization of \eqref{x_reciprocidad:eq:delta-sturm}.

An operator presentation makes memory explicit. On \(\ell^2(\ZZ)\), let \(T|n\rangle=|n+1\rangle\), \(V_9|n\rangle=\zeta_9^n|n\rangle\), \(W_\delta|n\rangle=e^{2\pi i n\delta}|n\rangle\) and \(D_M|n\rangle=\lfloor n/9\rfloor|n\rangle\). On finitely supported vectors,
\begin{equation}
 V_9T=\zeta_9TV_9,\qquad
 W_\delta T=e^{2\pi i\delta}TW_\delta,\qquad
 [D_M,T^9]=T^9.
 \label{x_reciprocidad:eq:weyl-relojes}
\end{equation}
The first two relations are Weyl covariances; the third expresses that a nonadic turn increases the memory degree. The base rotation is abelian, while the algebra of observables and translations is not. Inversion \(|n\rangle\mapsto|-n\rangle\) conjugates forward motion to backward motion. This realization provides an aperiodic helical suspension with finite phase return, not a periodic word of length nine.


% BEGIN OWNER /Users/ruben/Documents/New project/output/EXTREMA_MEDIA_RAZON_RECIPROCIDAD_ES_EN_20260918/ARTICULO/ES/source/figures/holonomia.tex
\begin{figure}[htbp]
\centering
\begin{tikzpicture}[x=1cm,y=1cm,>=Latex,font=\small]
\begin{scope}[xshift=1.8cm,yshift=.3cm]
\draw[->,gray] (0,-.5)--(0,3.65) node[above] {\(n/9\)};
\foreach \k in {0,1,2,3}{
 \draw[gray!60,dashed] (-1.1,\k)--(1.1,\k);
 \node[left,font=\scriptsize] at (-1.1,\k){\(\k\)};
}
\draw[black,line width=.65pt,domain=0:27,samples=260,smooth,variable=\n]
 plot ({cos(40*\n)},{\n/9+.14*sin(40*\n)});
\foreach \n in {0,...,27}{
 \fill ({cos(40*\n)},{\n/9+.14*sin(40*\n)}) circle (.026);
}
\foreach \k in {0,1,2,3}{\fill (1,\k) circle (.045);}
\node[align=center,text width=3.3cm,font=\scriptsize] at (0,-1.02)
 {Same phase in \(n=0,9,18,27\);\\different memory degree.};
\end{scope}
\begin{scope}[xshift=4.2cm,yshift=.5cm]
\node[anchor=west] at (0,3.1){\(z_n(0)=\lfloor(n+1)\delta\rfloor-\lfloor n\delta\rfloor\)};
\draw[->] (0,1.7)--(9,1.7) node[right]{\(n\)};
\foreach \n in {21,43,65,87,109,131,152,174,196,218}{
 \draw[line width=.6pt] ({\n*.039},1.7)--({\n*.039},2.25);
 \fill ({\n*.039},2.25) circle (.033);
 \node[below,font=\tiny] at ({\n*.039},1.7){\n};
}
\foreach \a/\b/\g in {21/43/22,43/65/22,65/87/22,87/109/22,109/131/22,131/152/21,152/174/22,174/196/22,196/218/22}{
 \node[font=\scriptsize] at ({(\a+\b)*.0195},2.55){\g};
}
\node[anchor=west,align=left,text width=8.3cm] at (0,.55)
 {Irrational encoding determines the intervals between events. Return of the nonadic phase does not turn this sequence into a periodic word.};
\end{scope}
\end{tikzpicture}
\caption{Complementary representations of transport. Left, projection of the helix \((\cos(2\pi n/9),\sin(2\pi n/9),n/9)\): advancing nine steps recovers phase and increases height. Right, the first ten events of the lower word for \(\delta=\log_{10}(10/9)\), with exact gaps \(21\) or \(22\). The helical representation introduces no additional physical metric.}
\label{x_reciprocidad:fig:holonomia-esturmiana}
\end{figure}

% END OWNER


For this count we fix the intercept \(\theta=0\) and accumulated endpoints \((N_0,\ldots,N_4)=(0,729,972,999,1000)\), corresponding to the cubic order \((729,243,27,1)\). The lower counts are \(\lfloor N_j\delta\rfloor-\lfloor N_{j-1}\delta\rfloor\), giving \((33,11,1,0)\). The upper counts are \(\lceil N_j\delta\rceil-\lceil N_{j-1}\delta\rceil\), with the same zero initial boundary, giving \((34,11,1,0)\). These values depend on the intercept and stratum order; they are not asserted for an arbitrary initial phase. Their totals are forty-five and forty-six. The increment corresponds to a single boundary unit, located in the first stratum of this ordered partition. The numbers thirty-four and eleven thus appear in the same count preceding any units chart; identifying them with physical exponents additionally requires the corresponding dimensional derivation.

The designation \emph{aperiodic crystal} is used here for this symbolic order, with specified complexity, spectrum and return. Zero topological entropy does not by itself equal zero physical dissipation: a thermodynamic realization requires defined energetic dynamics and an erasure operation. Likewise, this Sturmian reading does not replace all TPK operators or by itself prove periodicity or aperiodicity of the complete output of \(\alpha\). Its function is to make explicit an order structure that mere calendar return does not describe.


% BEGIN OWNER /Users/ruben/Documents/New project/output/EXTREMA_MEDIA_RAZON_RECIPROCIDAD_ES_EN_20260918/ARTICULO/ES/source/sections/revision_monodromia.tex
% Adición acumulativa. Insertar después de la suspensión esturmiana existente.
% Conservar íntegro el texto y la figura fig:holonomia-esturmiana.
\subsection{Nonadic connection, holonomy and monodromy representation}
\label{x_reciprocidad:mono-ampl:conexion}

The already constructed regional connection is now studied through three objects:
transport of the complete state, the representation of its turns, and the
suspension of its itineraries. The APP--TRIT--TPK operations and the
initial catalog retain their preceding definitions. The necessary
distinction here concerns what each representation retains of the
same history (Sections~\ref{x_reciprocidad:sec:extension} and~\ref{x_reciprocidad:sec:generacion}).

At prolongation index \(k\), a state chart preserves
\[
 \widehat x_k=(B_k,C_k,p_k,c_k,\Sigma_k,W_k,g_k,m_k).
\]
Here \(B_k\) combines the three regional blocks, \(C_k\) is the cylinder,
\(p_k\) the prefix, \(c_k\) the carry register, \(\Sigma_k\) the boundary
signature and \(W_k\) the incidence information. The phase is
\(g_k=1+[(k-1)]_9\); \(m_k\) counts complete turns. The reduction
\(\beta_k=[m_k]_{12}\) preserves the dodecaphasic position, while the
integer \(m_k\) distinguishes its successive turns.

\begin{definition}[Transport of a nonadic turn]
\label{x_reciprocidad:mono-ampl:def-vuelta}
Let \(\mathcal R_k\subseteq\widehat X_k\times\widehat X_{k+1}\)
be the admissible extension relations of TPK. The realization of
holonomy at level \(k\) is the relational composition
\[
 \Gamma_{9,k}=\mathcal R_{k+8}\circ\cdots\circ\mathcal R_k.
\]
For \(s\geq1\), its prolongation to \(s\) turns is
\[
 \Gamma_{9s,k}=
 \Gamma_{9,k+9(s-1)}\circ\cdots\circ\Gamma_{9,k}.
\]
\end{definition}

The relational formulation preserves admissible continuations. On
an individualized history, each arrow leads to the state actually
prolonged; the isolated visible block may still admit more than one
continuation. This difference is precisely the function of memory and
boundary conditions in determining the trajectory.

\begin{proposition}[Iteration of return and accumulated memory]
\label{x_reciprocidad:mono-ampl:vueltas}
In every admissible composition of \(s\) turns,
\[
 g_{k+9s}=g_k,\qquad m_{k+9s}=m_k+s,\qquad
 \beta_{k+9s}=\beta_k+s\pmod{12}.
\]
\end{proposition}
\begin{proof}
The one-turn rule preserves phase and increments the counter by one
unit. Applying it to the already transported state repeats exactly these
two operations. Induction gives the first two equalities; reducing the
second modulo \(12\) gives the third.
\end{proof}

The counter thus provides an integer witness to the difference between
turns. After \(12\) turns the pair of finite phases is restored, but
the counter increases by \(12\). The nonsplit extension
\eqref{x_reciprocidad:eq:extension108} expresses calendar composition through its
cocycle; the trajectory also preserves the integer lift of that datum.
Conservation here means transport of prior contributions,
not immobility of all coordinates.

\subsubsection{Bilateral representation of the return index}

The shift \(T\) of \eqref{x_reciprocidad:eq:weyl-relojes} admits an
expression in phase and turn coordinates. Number the phases from \(0\) to
\(8\), by shifting the origin relative to \(g_k\), and
write \(k=9m+r\), with \(0\leq r<9\), on the integer line
covering the phase cycle. On
\(\mathcal H_{\rm ind}=\ell^2(\mathbb Z)\), the basis
\(|k\rangle=|r,m\rangle\) expresses that same unit shift
\[
 T|r,m\rangle=
 \begin{cases}
 |r+1,m\rangle,&r<8,\\
 |0,m+1\rangle,&r=8.
 \end{cases}
\]
The inverse decreases \(r\) by one unit when \(r>0\), and takes
\(|0,m\rangle\) to \(|8,m-1\rangle\). Passing through the phase origin
therefore incorporates the exact quotient unit.

\Needspace{8\baselineskip}
\begin{proposition}[Representation of the cycle's monodromy]
\label{x_reciprocidad:mono-ampl:representacion}
Let \(S=T^9\). Then
\[
 S|r,m\rangle=|r,m+1\rangle,\qquad
 \mathbb Z\longrightarrow\mathcal U(\mathcal H_{\rm ind}),
 \quad s\longmapsto S^s
\]
is a faithful unitary representation of the cycle's turn group.
\end{proposition}
\begin{proof}
Nine shifts preserve the remainder of division by \(9\) and
increment the quotient. The powers satisfy
\(S^{s+t}=S^sS^t\), and \(S^s|r,m\rangle=|r,m+s\rangle\) distinguishes
every integer \(s\). Permutation of an orthonormal basis is unitary.
\end{proof}

This representation corresponds to the covering of the cycle, whose turn
group is \(\mathbb Z\), and represents phase with its counter. The
TPK state also contains the trajectory coordinates already described.
The bilateral extension admits negative indices as coordinates of the
covering; a trajectory starting from an initial condition uses
its admissible half-axis. An operator inverse of the index is thus
available without turning each extension relation into a bijection of the
complete state.

With \(D_M|k\rangle=\lfloor k/9\rfloor|k\rangle\), the identity
\([D_M,S]=S\) expresses the memory increment on the dense domain of
finitely supported vectors. If \(\mathcal I|k\rangle=|-k\rangle\) and
\(P_0\) projects onto indices divisible by \(9\), we also
obtain
\begin{equation}
 \mathcal I T\mathcal I=T^{-1},\qquad
 \mathcal I D_M\mathcal I=-D_M-(I-P_0).
 \label{x_reciprocidad:mono-ampl:inversion-contador}
\end{equation}
Indeed, for \(k=9m+r\),
\(\lfloor-k/9\rfloor=-m\) if \(r=0\) and equals \(-m-1\) if \(r>0\).
Inversion thus preserves the boundary term of Euclidean division.

\subsection{Block clock and capacity clock}
\label{x_reciprocidad:mono-ampl:relojes}

The index of a refinement map and the number of available decimal
positions have related but different laws. To
express them, reserve \(n\) for the ternary-block index and define
\[
 \rho=\log_{1000}729=\log_{10}9,\qquad
 \delta=1-\rho,\qquad
 \mathcal C_n=\lfloor(n+1)\rho\rfloor,\quad n\geq0.
\]
The letter \(\mathcal C_n\) here denotes capacity, not the dodecaphasic
register \(K\). At the canonical phase origin, the preceding mechanical word
provides
\[
 Z_n=\sum_{j=0}^{n}z_j(0)=\lfloor(n+1)\delta\rfloor.
\]

\begin{proposition}[Exact balance between capacity and vacancies]
\label{x_reciprocidad:mono-ampl:balance-capacidad}
For \(n\geq0\) and, in the second identity, \(n\geq1\),
\[
 \mathcal C_n=n-Z_n,\qquad
 \mathcal C_n-\mathcal C_{n-1}=1-z_n(0).
\]
\end{proposition}
\begin{proof}
Irrationality of \(\delta\) implies
\(\lceil(n+1)\delta\rceil=\lfloor(n+1)\delta\rfloor+1\).
Then
\[
 \lfloor(n+1)(1-\delta)\rfloor
 =n+1-\lceil(n+1)\delta\rceil=n-Z_n.
\]
Subtracting two consecutive identities completes the proof.
\end{proof}

An event \(z_n=1\) adds a block to the history while preserving capacity
\(\mathcal C_n\); an event \(z_n=0\) increments both counters. In
any interval, capacity increments and vacancies sum to
exactly its length. Reducing modulo \(q\), for \(q=9\) or
\(q=108\), gives
\[
 [\mathcal C_n]_q=[n]_q-[Z_n]_q.
\]
Uniform block phase and capacity phase are related through
integrated vacancy memory. In particular, capacity
retention and a complete turn of the connection have distinct
mechanisms, although both can preserve a phase observation.

Monotone inversion of the clock gives a first arrival time:
\[
 T_{\rm cap}(a)=\min\{n:\mathcal C_n=a\}
 =\left\lceil\frac a\rho\right\rceil-1,
 \qquad a\geq1.
\]
With \(\sigma=\rho^{-1}-1\), the time spent gaining \(b\) units
of capacity is
\[
 T_{\rm cap}(a+b)-T_{\rm cap}(a)
 =b+\lceil(a+b)\sigma\rceil-\lceil a\sigma\rceil.
\]
The difference of ceilings takes one of the two integers adjacent to
\(b\sigma\). This yields \(9\) or \(10\) blocks to gain
\(9\) units of capacity, and \(113\) or \(114\) to gain \(108\).
These are times of the inverse clock, measured from a first arrival; the
turn counter is always interpreted in the index to which it belongs.

\subsection{Symbolic suspension and aperiodic time crystal}
\label{x_reciprocidad:mono-ampl:cristal}

Rotation encoding combines phase periodicity, local recurrence
and global aperiodicity. Let \(\mathcal X_\delta\subseteq\{0,1\}^{\mathbb Z}\)
be the closure, in the product topology, of the translates of the
bilateral word \(z_n(\theta)\). Denote by \(\sigma_{\rm sh}\) the
shift of its sequences. Retaining a nonadic phase, the step is
\[
 F(r,z)=([r+1]_9,\sigma_{\rm sh}z),\qquad
 (r,z)\in C_9\times\mathcal X_\delta.
\]

\begin{definition}[Symbolic aperiodic time crystal]
\label{x_reciprocidad:mono-ampl:def-cristal}
The unit-roof suspension of the preceding system is the space
\[
 \mathcal S_\delta=
 \bigl((C_9\times\mathcal X_\delta)\times\mathbb R\bigr)/\sim,
 \qquad (y,t+1)\sim(Fy,t),
\]
with flow \(\Phi_s[y,t]=[y,t+s]\). This temporal-order model, together with its
vacancy encoding and turn-counter representation, is called a \emph{symbolic
aperiodic time crystal}.
\end{definition}

The adjective temporal designates the action of the flow and of its return section;
aperiodicity characterizes its itineraries. The frequency of ones in each
sequence is \(\delta\): the uniform discrepancy bound of less than one,
obtained by telescoping in any window, persists upon taking the closure.
A periodic sequence would have a rational frequency. Therefore,
\(\mathcal X_\delta\) contains no periodic orbits. The suspension
also has no closed orbits: a return of the flow would require an integer
time and periodicity of the itinerary.

Finite words do recur. The itineraries of length \(m\)
arise from a partition of the circle by \(m+1\) points; each
interval is visited recurrently by the irrational rotation. Even
when times are restricted to a fixed phase, the step \(9\delta\) remains
irrational. This is why each word appears in the
\(9\) phases and why the complexities already established are
\[
 p(m)=m+1,\qquad p_9(m)=9(m+1),\qquad p_3(m)=3(m+1).
\]
The recurrence of local configurations coexists with the absence of a
global period. Its topological entropy is zero because these complexities
grow linearly.

The realization by the observables \(V_9,W_\delta\) and the shift
\(T\) expresses the same order through the covariances
\eqref{x_reciprocidad:eq:weyl-relojes}. The frequencies of the uniform clock belong to
\(\frac19\mathbb Z+\delta\mathbb Z\), modulo \(1\). The capacity
clock additionally preserves the relation
\([\mathcal C_n]_9=[n]_9-[Z_n]_9\); it is not replaced by the uniform phase
when interpreting the returns. The term \emph{time crystal} is thus
defined through a mathematical dynamical system. Its physical realization
requires specifying an energy dynamics, its temporal observables
and its stability, data distinct from the present symbolic construction.

\subsection{Mirror symmetry of the regions and dual realization}
\label{x_reciprocidad:mono-ampl:regiones}

The mirror interchange acts on the regional block
\(B=(u^\pi,u^e,u^\varphi)\in(\mathbb F_3^6)^3\) by
\[
 \mathfrak cB=(u^e,u^\pi,u^\varphi).
\]
Its fixed axis and its aggregate are, respectively, \(u^\varphi\) and
\(s=u^\pi+u^e\). This fixing refers to the involution in a fiber;
both coordinates may evolve during nonadic transport.
With \(q=u^\pi+u^e+u^\varphi\),
\(a=u^\pi+u^e-u^\varphi\) and \(d=u^\pi-u^e\), one recovers
\[
 u^\varphi=2(q-a),\qquad s=q-u^\varphi,\qquad
 u^\pi=2(s+d),\qquad u^e=2(s-d).
\]
The equalities hold componentwise in \(\mathbb F_3\), where
\(2^{-1}=2\). The mirror preserves \(q,a,s\) and reverses \(d\); this
antisymmetric coordinate singles out the distribution that the aggregate preserves
without distinguishing it. The reversal \(\mathcal I\) of the temporal index and the
involution \(\mathfrak c\) of channels therefore act on
different data.

The phase return has an explicit witness:
\[
 \begin{aligned}
 B_1&=(012222\mid121221\mid112202),\\
 B_{10}&=(101222\mid112221\mid112002).
 \end{aligned}
\]
Both positions have phase \(1\); the block, the axis and the memory have
changed. Likewise, synchronization of the censuses in phases \(8\) and
\(9\) cancels their relative differences, while the diagonal component
changes from \((59,59,59)\) to \((43,43,43)\). Equality of a relative
observable does not exhaust the state that determines it.

The dual realization preserves that architecture. In section
\ref{x_reciprocidad:exc:doble-lectura}, the same prefix determines nested
cylinders on the one hand and a sequence of encoded words with their frames on the other.
The first map passes to the limit by contraction of diameters; the second
satisfies \(r_{n,m}Q_n=Q_mp_{n,m}\) and defines \(Q_\infty\).
Figure \ref{x_reciprocidad:mono-ampl:fig-regional} anticipates these two maps.
The subsequent incidence path constructs codes, lattices and the
realization of Moonshine in their own domains; the automorphism group
acts on that realization, not on the digits through an implicit
identification.


% BEGIN OWNER /Users/ruben/Documents/New project/output/EXTREMA_MEDIA_RAZON_RECIPROCIDAD_ES_EN_20260918/ARTICULO/ES/source/figures/monodromia_regional_ampliacion.tex
% Figura acumulativa; destino figures/monodromia_regional_ampliacion.tex.
\begin{figure}[p]
\centering
\begin{tikzpicture}[x=1cm,y=1cm,>=Latex,font=\small,
  flow/.style={->,line width=.55pt},
  ann/.style={align=center,font=\scriptsize}]
\node[font=\bfseries] at (6.9,9.2)
  {Nonadic connection and regional coordinates};
\begin{scope}[shift={(3.4,6.8)}]
 \foreach \k in {1,...,9}{
  \pgfmathsetmacro{\ang}{90-40*(\k-1)}
  \coordinate (p\k) at (\ang:1.52);
  \fill (p\k) circle (1.4pt);
  \node[fill=white,inner sep=1.1pt,font=\scriptsize]
    at (\ang:1.83) {\(\k\)};
 }
 \foreach \k in {1,...,8}{
  \pgfmathtruncatemacro{\next}{\k+1}
  \draw[flow,shorten >=3pt,shorten <=3pt] (p\k)--(p\next);
 }
 \draw[flow,shorten >=3pt,shorten <=3pt] (p9)--(p1);
 \node[ann] at (0,.1) {\(g_k\in C_9\)\\\(\Gamma_{9,k}\)};
 \node[ann] at (0,-2.23) {One turn: \(g\mapsto g\), \(m\mapsto m+1\)};
\end{scope}
\node[ann,text width=5.25cm] at (10.1,7.75)
 {\(B_k=(u_k^\pi,u_k^e,u_k^\varphi)\)\\
 regional block in a fiber};
\node at (8.55,6.5) {\(u^\pi\)};
\node at (11.65,6.5) {\(u^e\)};
\draw[<->,line width=.65pt] (8.95,6.5)--node[above]{\(\mathfrak c\)}(11.25,6.5);
\draw[densely dashed,gray] (10.1,5.72)--(10.1,7.23);
\node[fill=white,inner sep=2pt] at (10.1,5.6){\(u^\varphi\)};
\node[ann] at (10.1,4.6)
 {Axis \(u^\varphi\) and aggregate \(u^\pi+u^e\)\\
 invariant under \(\mathfrak c\)};
\draw[flow] (5.45,6.8)--(7.25,6.8);
\draw[gray!60,line width=.3pt] (.55,4.02)--(13.25,4.02);
\node[ann,text width=12cm] (hist) at (6.9,3.47)
 {Compatible history: prefixes, cylinders, orientation, quotients,
 incidence frames and memory};
\coordinate (fork) at (6.9,2.85);
\draw[flow] (hist.south)--(fork);
\draw[flow] (fork)--(3.6,2.2);
\draw[flow] (fork)--(10.3,2.2);
\node[ann,text width=5.6cm] at (3.6,1.65)
 {Archimedean realization\\
 \(I_{n+1}\subseteq I_n,\quad |I_n|=729^{-n}\)\\
 limit of the regional prefixes};
\node[ann,text width=5.6cm] at (10.3,1.65)
 {Incidence realization\\
 \(r_{n,m}Q_n=Q_mp_{n,m}\)\\
 encoded words and compatible frames};
\node[ann,text width=5.6cm] at (3.6,.42)
 {\(\pi,\varphi,e\); dodecaphasic determination of \(\alpha\)};
\node[ann,text width=5.6cm] at (10.3,.42)
 {Witt--Golay--Leech--\(V^\natural\)\\
 automorphisms of the exceptional realization};
\draw[flow] (3.6,1.03)--(3.6,.78);
\draw[flow] (10.3,1.03)--(10.3,.78);
\end{tikzpicture}
\caption{Nine phases of a connection and two realizations of the compatible
history. The cycle represents the phase; each turn increments the counter.
The regional diagram represents the involution \(\mathfrak c\), which
interchanges closure and propagation and fixes the autoscale axis in a fiber.
This invariance does not require the axis to be constant during transport.
The two lower maps use the same prefix: one determines
its numerical limit and the other preserves words and incidence frames. The
maps of this second realization are specified in section
\ref{x_reciprocidad:exc:doble-lectura}; the exceptional path is developed afterward.}
\label{x_reciprocidad:mono-ampl:fig-regional}
\end{figure}

% END OWNER



% BEGIN OWNER /Users/ruben/Documents/New project/output/EXTREMA_MEDIA_RAZON_RECIPROCIDAD_ES_EN_20260918/ARTICULO/ES/source/sections/recta_nonadica_recuperada.tex
\clearpage
\subsubsection{Rectilinear representation of the phases and the return}
\label{x_reciprocidad:mono-recta:desarrollo}

The representation on the segment \([0,10]\) preserves the positional
origin of APP and orders its nine interior marks. In this diagram,
the integers \(1,\ldots,9\) are phase representatives. The horizontal
position indicates the order of the charts; the upper labels
specify regional operations. Figure
\ref{x_reciprocidad:mono-recta:figura} thus brings together the initial line, the fixing of the
autoscale axis, the mirror distribution of the two channels and the
return with memory (section~\ref{x_reciprocidad:mono-ampl:conexion}).


% BEGIN OWNER /Users/ruben/Documents/New project/output/EXTREMA_MEDIA_RAZON_RECIPROCIDAD_ES_EN_20260918/ARTICULO/ES/source/figures/recta_nonadica_0_10.tex
% Recuperación de fig_nonadic_route_mini.tex de la síntesis académica.
% Los rótulos designan operaciones regionales; la suma se efectúa en F_3^6.
\begin{figure}[H]
\centering
\begin{tikzpicture}[x=1.12cm,y=1cm,font=\small,>=Latex,
  guide/.style={line width=.45pt,black!55}]
 \draw[line width=.8pt] (0,0)--(10,0);
 \foreach \k in {0,...,10}{
  \draw[guide] (\k,-.09)--(\k,.09);
  \node[below=2.5mm] at (\k,0) {$\k$};
 }
 \foreach \k in {1,...,9}{\fill (\k,0) circle (1.7pt);}
 \node[anchor=west,font=\footnotesize] at (0,3.05)
   {Generating segment $[0,10]$ and nine phase representatives};
 \draw[guide] (5,.12)--(5,1.7);
 \node[above,align=center,font=\footnotesize] at (5,1.7)
   {fixing of the axis\\$u^\varphi$};
 \draw[decorate,decoration={brace,amplitude=4pt},guide]
   (6,.65)--(8,.65);
 \node[above=5pt,align=center,font=\footnotesize] at (7,.65)
   {regional aggregate\\$u^e+u^\pi$};
 \draw[guide] (9,.12)--(9,1.7);
 \node[above,align=center,font=\footnotesize] at (9,1.7)
   {mirror involution\\$\pi\leftrightarrow e$};
 \draw[->,line width=.7pt] (9,-.7)
   .. controls (9,-1.55) and (1,-1.55) .. (1,-.7);
 \node[fill=white,inner sep=3pt,font=\footnotesize] at (5,-1.23)
   {return $9\to1$;\quad $m\mapsto m+1$};
 \node[font=\footnotesize] at (5,-1.95)
   {$w_6\to w_{12}\to w_{18}\to w_{24}\to w_{30}\to R_{36}\to G_9$};
\end{tikzpicture}
\caption{Rectilinear representation of the nonadic connection. The mark
\(5\) locates the first strong saturation and the regional
autoscale axis; the brace \(6\)--\(8\) identifies the ternary aggregate
fixed in each fiber; phase \(9\) selects the closure
orientation. The return preserves the phase and increments the memory. The
abscissae number operational positions, while the labels
\(u^\pi,u^e,u^\varphi\) designate blocks of \(\mathbb F_3^6\).}
\label{x_reciprocidad:mono-recta:figura}
\end{figure}

% END OWNER


The fixing of the axis is formulated in the fiber of a given boundary. If
\(q=u^\pi+u^e+u^\varphi\) and
\(a=u^\pi+u^e-u^\varphi\), then
\[
 u^\varphi=2(q-a),\qquad s=u^\pi+u^e=q-u^\varphi
 \quad\text{in }\mathbb F_3^6.
\]
The data \((q,a)\) determine the axis and the aggregate; the ordered distribution
of \(s\) preserves the mirror freedom resolved by the prolongation.
The fifth phase has a local solution and an extension, with boundary
datum \(c=111111\) and residual signature \(\rho_W=000\). This is the
meaning of its first strong saturation. The symbol \(\varphi\)
placed above the mark \(5\) identifies the regional axis involved
in that operation.

The intermediate phases preserve this decomposition within each fiber,
while their boundary data evolve. The complete census of
local multiplicities and extensions of the distinguished branch is
\[
\begin{array}{c|rrrrrrrrr}
\text{phase}&1&2&3&4&5&6&7&8&9\\\hline
N_{\rm loc}&3&2&5&4&1&1&3&3&2\\
N_{\rm ext}&2&5&4&1&1&3&3&2&1
\end{array}
\]
Phase \(6\) combines local uniqueness and three prolongations; phase
\(7\) resolves a triple mirror fiber; phase \(8\) synchronizes
the three regional censuses. Phase \(9\) preserves two local
distributions with the same axis and the same aggregate, one of which
extends. These counts express different functions within a
single connection, and justify the marks on the line.

In particular, for phases \(6,7,8\) the regional sums are
respectively \(012100\), \(101001\) and \(112000\). The conservation
represented by the brace is an invariance with respect to the mirror
distribution in each fiber. The subsequent real evaluation applies the
readers of the compatible cylinders to the complete histories; the
internal sum represented in the figure belongs to the ternary domain.

Finally, the arc \(9\to1\) represents the passage from one turn to the
next. The formulas in section
\ref{x_reciprocidad:mono-ampl:conexion} preserve the phase remainder and increment the
turn counter. The rectilinear diagram and the cyclic diagram in
Figure \ref{x_reciprocidad:mono-ampl:fig-regional} therefore express two complementary
aspects of the same holonomy: the order of the positions and the
composition of the return on the state with memory.

% END OWNER


\subsection{Temporal transport and tritic exponential composition}
\label{x_reciprocidad:mono-ampl:euler}

The Euler relations in section \ref{x_reciprocidad:euler-trit-articulacion}
describe the local realizations of the three regimes. For a fixed
signature \(\tau\), the generator satisfies \(J_\tau^2=-\tau I\); algebraic
conjugation takes \(J_\tau\) to \(-J_\tau\). Consequently,
\[
 \exp(sJ_\tau)\exp(tJ_\tau)=\exp((s+t)J_\tau),\qquad
 \overline{\exp(tJ_\tau)}=\exp(-tJ_\tau),
\]
and the product of an evolution with its conjugate is unity. In the elliptic
type it is realized through \(\cos^2t+\sin^2t=1\); in the hyperbolic type,
through \(\cosh^2t-\sinh^2t=1\); in the parabolic type,
\((I+tN)(I-tN)=I\), since \(N^2=0\). The same conservation law
admits three quadratic expressions.

The elliptic realization closes a turn; the nilpotent one preserves a
linear displacement; the hyperbolic one separates expansion and contraction with
unit determinant. In the regional parameters already generated,
\[
 \exp(\pi J_{\mathrm e})=-I,\qquad
 \exp(J_{\mathrm p})=I+J_{\mathrm p},\qquad
 \exp(2\log\varphi\,J_{\mathrm h})=F_{\rm av}^{\,2}.
\]
The evaluation \(\exp(I)=eI\) recognizes the unity of the common exponential
law. The number \(e\) is not assigned exclusively to the parabolic regime.

Transporting a local realization through an isomorphism \(U\) preserves
these relations: if \(J'=UJU^{-1}\), the exponential series gives
\(U\exp(tJ)=\exp(tJ')U\). The identity transports the regime and its
norm. A transition between different quadratic types instead belongs
to the selection of sheets and regimes of TPK; it is not a conjugation
that changes the value of \(J^2\). The temporal product therefore preserves the
order of the transitions and the information of their endpoints.

This articulation distinguishes the closure of a local representation and the
holonomy of a history. It is possible to have \(\exp(2\pi J_{\mathrm e})=I\)
at the same time as \(S|r,m\rangle=|r,m+1\rangle\). The first equality
restores the orientation of the regime; the second records another turn. The
discrete structure preserves both: the local return and the accumulated
provenance of the state that realizes it.

% END OWNER



% BEGIN OWNER /Users/ruben/Documents/New project/output/EXTREMA_MEDIA_RAZON_RECIPROCIDAD_ES_EN_20260918/ARTICULO/ES/source/sections/primos_retornos.tex
% Recuperación del producto nativo y su selector, no incorporación de RH.
\subsection{Multiplicative irreducibility and return periods}
\label{x_reciprocidad:primos-retornos}

The distinction between structure, evaluation and phase also appears in the
multiplicative sector. The local product of APP contains both the residue
and the quotient; its iteration allows normal forms and
factorizations to be constructed before examining their decimal periods. This development
places the expression \emph{prime mode} in a precise domain and relates
irreducibility to the operations already defined, not to an isolated
coincidence of periods (section~\ref{x_reciprocidad:primos-retornos}).

Let
\[
 \mathcal W_9=\{\varnothing\}\sqcup
 \bigsqcup_{\ell\geq1}\{1,\ldots,9\}^{\ell}.
\]
Sequences are ordered from the least significant digit. The evaluation
\(\operatorname{ev}_9(u)=\sum_j u_j9^j\), with
\(\operatorname{ev}_9(\varnothing)=0\), is a bijection onto the nonnegative
integers: for a positive integer, dividing \(n-1\) by \(9\)
uniquely determines the first digit and the continuation. This is
bijective nonadic numeration; it preserves the boundary representative
\(9\).

The sum \(\boxplus_\#\) is obtained by applying the additive APP sheet to the
digits at each position and propagating their quotients. If there is only one digit,
it constitutes the provisional residue; a nonzero incoming quotient is
normalized through a second additive operation. The sum of the
quotients is transported to the next position. To multiply a
sequence by a digit \(d\), the multiplicative cell produces
\(u_jd=9q_j+r_j\), and the incoming quotient \(c_j\) is incorporated through
\[
 (w_j,c_{j+1})=
 \begin{cases}
 (r_j,q_j),&c_j=0,\\
 \bigl(R^+(r_j,c_j),q_j+q^+(r_j,c_j)\bigr),&c_j>0.
 \end{cases}
\]
The procedure begins with \(c_0=0\) and incorporates the final quotient when it is
nonzero. During addition, \(c_j\leq2\); during multiplication by a digit,
\(c_j\leq9\). The local operations thus receive arguments within
their domain. Denote this second procedure by \(d\odot_\#u\),
with \(d\odot_\#\varnothing=\varnothing\).

For \(v=(v_0,\ldots,v_{m-1})\), the full product is constructed by
the recurrence
\[
 a_m=\varnothing,\qquad
 a_j=(9\odot_\#a_{j+1})\boxplus_\#(v_j\odot_\#u),
 \qquad u\boxtimes_\#v=a_0.
\]
The rule composes the two local sheets. Its evaluation satisfies
\begin{equation}
 \operatorname{ev}_9(u\boxtimes_\#v)
   =\operatorname{ev}_9(u)\operatorname{ev}_9(v).
 \label{x_reciprocidad:eq:primos-entrelazamiento}
\end{equation}
Indeed, each transition preserves the partial sum plus the pending
quotient weighted by the power of \(9\). The recurrence then gives
\(\operatorname{ev}_9(a_j)=9\operatorname{ev}_9(a_{j+1})+
v_j\operatorname{ev}_9(u)\), whose iteration proves
\eqref{x_reciprocidad:eq:primos-entrelazamiento}. Uniqueness of the normal form transports
associativity and commutativity to the constructed product; its unit is \((1)\).

\begin{definition}[Reduced multiplicative coproduct]
On the vector space of finite support generated by the nonempty
normal forms, define
\[
 \widetilde\Delta_\# e_w
 =\sum_{\substack{u\boxtimes_\#v=w\\u,v\ne(1)}}e_u\otimes e_v.
\]
A nonunit normal form is irreducible when its image under
\(\widetilde\Delta_\#\) is zero.
\end{definition}

The factorization fibres are finite by
\eqref{x_reciprocidad:eq:primos-entrelazamiento}. In the completion
\(\mathcal H_\#=\ell^2(\mathcal W_9\setminus\{\varnothing\})\),
the images of distinct forms have disjoint supports. If
\(d_\#(w)\) counts the ordered nonunit factorizations, the domain
of the closure is
\[
 \operatorname{Dom}(\overline{\widetilde\Delta_\#})
 =\left\{x\in\mathcal H_\#:
       \sum_w d_\#(w)|x_w|^2<\infty\right\}.
\]
Indeed, the squared norm of the image is that weighted sum. Therefore,
the operator
\[
 A_\#=(\overline{\widetilde\Delta_\#})^*
                  \overline{\widetilde\Delta_\#}
\]
is diagonal, positive, and self-adjoint: its eigenvalue at \(e_w\) is the
number of ordered nonunit factorizations of \(w\). The projection
\[
 P_{\rm irr}=(I-|e_{(1)}\rangle\langle e_{(1)}|)
                    \mathbf1_{\{0\}}(A_\#)
\]
selects exactly the irreducible forms. Through multiplicative
evaluation, these correspond to the ordinary primes. Selection
depends on the factorization structure and keeps the unit
state separate.

For a value \(n\geq2\) coprime to \(10\), decimal division induces the
residue permutation \(r\mapsto10r\pmod n\). Return of the residue
\(1\) has period
\[
 \tau_n=\operatorname{ord}_n(10).
\]
If \(n=\prod_p p^{a_p}\), the Chinese remainder theorem gives
\[
 \tau_n=\operatorname{mcm}_{p^{a_p}\parallel n}
                    \operatorname{ord}_{p^{a_p}}(10).
\]
The composite return synchronizes the returns of the prime powers.
For a prime \(p\notin\{2,3,5\}\), we have
\[
 \tau_p\mid p-1,\qquad
 M_p=\frac{10^{\tau_p}-1}{p}\in\mathbb N,
 \qquad M_p\equiv0\pmod9.
\]
The last congruence expresses the nonadic closure of the decimal repetition:
\(10^{\tau_p}-1\) is a multiple of \(9\), and \(p\) is invertible modulo
\(9\). The same congruence holds for any \(n\) coprime to
\(90\). In particular, \(7,13,77,91,143\) have decimal period \(6\),
although only the first two are irreducible. The coproduct and the period
therefore record different properties of the same constructed value.

This distinction makes precise the relationship with aperiodic refinement.
The vacancy sequence determines the increments of the capacity
index; factorization determines the multiplicative modes;
decimal division determines their return times. The nonadic
representative preserves phase closure, while the quotient and
factorization preserve information that this closure does not distinguish. The
link with the spectral constructions of the treatise resides in those
operators and their intertwinings. The proofs of their subsequent
analytic results retain their own development; they are not replaced by
the periodicity of a residual observation.

\subsubsection{Vacancy discrepancy and Dirichlet readouts}

The link also admits an exact analytic expression. The same
encoding \(z_n=z_n(0)\), for \(n\geq1\), is preserved, and
two subsequent observables are considered: the series over all indices and the
product over the indices selected by irreducibility. For
\(\operatorname{Re}s>1\),
\[
 Z_{\rm vac}^{+}(s)=\sum_{n\geq1}\frac{z_n}{n^s},\qquad
 Z_{\rm vac}^{\times}(s)=\prod_p(1-p^{-s})^{-z_p}.
\]
Both converge absolutely in that half-plane. Their common factor is the
density \(\delta\), but their discrepancies differ.

\Needspace{21\baselineskip}
\begin{proposition}[Separation of vacancy density]
We have
\[
 Z_{\rm vac}^{+}(s)=\delta\zeta(s)+H_{\rm vac}(s),
 \qquad H_{\rm vac}\in\mathcal O(\operatorname{Re}s>0).
\]
In \(\operatorname{Re}s>1\), with logarithms defined by the absolutely
convergent Euler series,
\[
 Z_{\rm vac}^{\times}(s)=\zeta(s)^\delta G_{\rm vac}(s),\qquad
 \log G_{\rm vac}(s)=
 \sum_p(z_p-\delta)\log(1-p^{-s})^{-1}.
\]
\end{proposition}
\begin{proof}
The centred partial sums satisfy
\[
 A_N:=\sum_{n=1}^N(z_n-\delta)
 =\lfloor(N+1)\delta\rfloor-N\delta,\qquad |A_N|<1.
\]
Abel summation represents the centred term by
\(s\int_1^\infty A_{\lfloor x\rfloor}x^{-s-1}\,dx\). The integral
converges locally uniformly in \(\operatorname{Re}s>0\),
proving holomorphy. The second identity follows by separating
\(z_p=\delta+(z_p-\delta)\) in the logarithm of the product, within
its domain of absolute convergence.
\end{proof}

The decomposition makes visible what the readout on primes adds: it examines
the same vacancy sequence on a subset determined by
factorization. Control of the discrepancy over all integers gives
the first holomorphic extension; the second observable preserves a
discrepancy of its own on primes. These are two applications of the same
discrete structure, with explicit analytic domains. This articulation
explains their relevance to the arithmetic core without transferring to this
article the proofs of the terminal results of the treatise.

% END OWNER

